% !TeX root = main.tex

The RCP solver is (in principle) complete for %orderable fragment~\cite{???}, %which subsumes
the straight-line fragment~\cite{ostrich} and chain-free fragment~\cite{chain19} of string constraints,
in the sense that the underlying proof system can show unsatisfiability of unsatisfiable instances from these fragments,
when applying the rules in an appropriate (easy) order.
However, the usage of priorities means the rules may be applied in a different order, causing the proof to fail.
In practice, failures for these fragments are rare.

In the straight-line fragment we require that the constraints are in normal form as in Section~\ref{sec:inference-rules},
and variables can be ordered $x_1, \ldots, x_n$
such that for $x_i$ there is exactly one equational constraint $x_i = f_i(x_1, \ldots, x_{i-1})$,
where $f_{i}$ is a string function such that the pre-image of a regular set is regular.
The definition of the chain-free fragment is a bit involved, in essence it relaxes the assumption on the form of the equation and the number of equations,
but still requires that there are no ``cycles'' of dependencies.
For the straight-line fragment, we can propagate backwards and the instance is unsatisfiable if and only if all branches are closed using Close rule.
For the chain-free fragment~\cite{chain19} we can adapt the original proof (which works for a different solver) to the rules used in OSTRICH2.

The CE-Str solver also provides a completeness guarantee for a version of the straight-line fragment that includes string constraints and linear integer arithmetic~\cite{atva2020}.
The supported string functions are those where the pre-image of cost-enriched automata constraints can be expressed using cost-enriched automata.
This includes $\sindexof$ and $\slength$ as well as some versions of $\sreplace$ and $\sreplaceall$.


%The orderable fragment is defined procedurally defined: we assume that the constraints are in normal form as in Section~\ref{sec:inference-rules}.
%In each round we mark variables that occur once (in all constraints) then, if there is a constraint $x = f(\overline y)$ such that
%each variable in $\overline y$ is marked and an image of regular set by $f$ is regular (i.e.\ for each regualr $L$ the $f(L)$ is also regular) then
%we remove the constrain (and continue);
%similarly, if $x$ is marked and the pre-image of $f$ of each regular set is a recognizable relation (i.e.\ it is of a form $\bigcup^k_{i=1}
%L^i_1\times \dots \times L^i_n$ for some regular sets $\{L_j^i\}$)
%then we remove the constraint.
%In particular, this algorithm is effective and yields an order in which to apply the \texttt{[Fwd-Prop]} and \texttt{[Bwd-Prop]} rules.

